\section{Multiscale stopping and all fixed marked moments}
\label{sec:multiscale-stopping}
A single clock window assigns the same exceptional-set cost to every
larger deviation. Summing instead over the actual deviation scales
recovers a fourth-root comparison. This estimate is uniform in the
location of a prescribed return mark; it does not take a maximum over
all marks inside the expectation. The stronger comparison will also
remove the loss of central rate on the wider Fourier band.

Retain the notation of \eqref{eq:moment-stopped-pair}:
\[
 Z_{n,k,R}=\mathsf H_{N^-_{k,R},N^+_{n-k,R},R},\qquad
 W_{n,k,R}=\mathsf H_{m_k,m_{n-k},R},\qquad
 m_j=\lfloor j/c_*\rfloor.
\]
In particular, recentering the orbit at its actual marked return gives
$Z_{n,k,R}(y)=U_{n,R}((F_R^*)^{-k}y)$, where
$U_{n,R}=J_{n,R}-n\bar G_R$. All norms in the stopping argument are
under the normalized section probability $\nu_R^*$. Constants may
depend on a fixed moment order, but not on $R,n,k$.

\subsection{Deviation shells rather than one exceptional set}
\begin{lemma}[Two bounds for the total clock deviation]
\label{lem:total-clock-deviation}
Set
\[
 D_{n,k,R}=|N^-_{k,R}-m_k|+|N^+_{n-k,R}-m_{n-k}|.
\]
For every fixed integer $p\ge2$ there are constants $C_p,a_0,c_0>0$
such that, for $n\ge2$ and $b\ge2$,
\begin{align}
 \nu_R^*(D_{n,k,R}>b)&\le C_p\frac{(n+b)^p}{b^{2p}},
       \label{eq:total-clock-polynomial}\\
 \nu_R^*(D_{n,k,R}>b)&\le C_pe^{a_0n-c_0b}.
       \label{eq:total-clock-exponential}
\end{align}
The second bound may exceed one. It will only be used beyond a fixed
multiple of $n$.
\end{lemma}
\begin{proof}
A deviation greater than $b$ requires one of the two clocks to deviate
by more than $b/2$. Equation~\eqref{eq:higher-moment-clock} and a
union bound prove \eqref{eq:total-clock-polynomial}, with a change of
constant for bounded $b$. This argument includes $k=0,n$; a zero-return
clock is identically zero.

Put $Q=N^-_{k,R}+N^+_{n-k,R}$. Under $\nu_R^*$, a backward clock of
$j$ returns has the same law as a forward clock of $j$ returns, since
$N^-_{j,R}((F_R^*)^jx)=N^+_{j,R}(x)$ and the actual return map preserves
that probability. The cumulative-return estimate therefore gives
\[
 \nu_R^*(Q>L)\le C e^{an-cL/2}\qquad(L\ge0)
\]
by a union bound, with no independence of the two clocks. Since
$D_{n,k,R}\le Q+m_k+m_{n-k}\le Q+c_*^{-1}n$, translating the threshold
proves \eqref{eq:total-clock-exponential}. For negative translated
thresholds, increase $a_0$ and $C_p$ so that the displayed right side
is at least one.
\end{proof}

\begin{theorem}[Fourth-root stopping in every fixed finite moment]
\label{thm:multiscale-marked-stopping}
For every fixed real $q\ge1$ there is $C_q<\infty$ such that
\begin{align}
 \|Z_{n,k,R}\|_q+\|W_{n,k,R}\|_q&\le C_q\sqrt n,
       \label{eq:all-stopped-moment-bound}\\
 \|Z_{n,k,R}-W_{n,k,R}\|_q&\le C_q n^{1/4}
       \label{eq:quarter-stopping}
\end{align}
for all $n\ge2$, $0\le k\le n$, and $R\in I$. Consequently, for every
bounded insertion $a$ supported on the section and every real $v$,
\begin{align}
 &\left|C^{[k],a}_{n,R}(v)
   -\int_M a(y)e^{iv\cdot W_{n,k,R}(y)/\sqrt n}\dd\nu(y)\right|
       \le CM_a|v|n^{-1/4}.
       \label{eq:quarter-characteristic-stopping}
\end{align}
For $q=2$, the proof of \eqref{eq:quarter-stopping} uses only the
fourth collision moments, the fourth-moment clock window, and the
cumulative-return tail.
\end{theorem}
\begin{proof}
Choose a fixed integer $p\ge2$ such that $2p>3q/2$. The deterministic
interval has length at most $c_*^{-1}n$. Its $q$th moment is bounded
by Lemma~\ref{lem:fixed-even-maximal} at a sufficiently large fixed
even order, followed by domination of the initial density.
For $Z$, split at $Q\le An$, with $Q$ as in the preceding proof and
$A$ a sufficiently large fixed constant. On this event, $|Z|$ is at
most the sum of two deterministic collision maxima of length
$\lceil An\rceil$. Their $q$th moments are $O_q(n^{q/2})$. On the
complement, boundedness of $h_R$ gives $|Z|\le C Q$. Integration of
the cumulative exponential tail bounds the remaining $q$th moment by
$C_q(1+n)^q e^{-c'n}$. This proves
\eqref{eq:all-stopped-moment-bound}.

We give the full multiscale argument for the difference. For each
$L\ge1$, let $B_L$ be the sum of the four prefix or suffix maxima
on the deterministic length-$\lceil L\rceil$ windows adjacent to
the two endpoints $-m_k$ and $m_{n-k}$. All these windows are permitted
to cross collision time zero. They are fixed intervals, not intervals
conditioned on the clocks. By Lemma~\ref{lem:fixed-even-maximal},
\begin{equation}\label{eq:deterministic-shell-maximal}
 \int B_L^{2p}\dd\nu_R^*\le C_p L^p.
\end{equation}
On $D_{n,k,R}\le L$ one has the pathwise bound
$|Z-W|\le B_L$. Let $b_0=\lceil\sqrt n\rceil$ and
$L_j=2^j b_0$ for $j\ge1$. Partition the sample space into
$\{D\le b_0\}$ and the disjoint sets
$\{L_j/2<D\le L_j\}$, $j\ge1$. The clocks are finite almost
everywhere, so these sets exhaust the section up to a null set.
The first set contributes at most $C_q b_0^{q/2}$ to
$\int|Z-W|^q\dd\nu_R^*$.

On every other shell, H\"older's inequality gives
\begin{equation}\label{eq:dyadic-shell-holder}
 \int_{\{L_j/2<D\le L_j\}}|Z-W|^q\dd\nu_R^*
 \le C_{p,q}L_j^{q/2}
       \nu_R^*(D>L_j/2)^{1-q/(2p)}.
\end{equation}
The exponent is positive. Fix $A$ large enough that the exponential
bound \eqref{eq:total-clock-exponential} at threshold $L_j/2$ is
at most $Ce^{-c_0L_j/4}$ whenever $L_j>An$.
For the shells $L_j\le An$, use
\eqref{eq:total-clock-polynomial} in
\eqref{eq:dyadic-shell-holder}. Their total is at most
\begin{align}
 C_{p,q,A}\sum_{L_j\le An} n^{p-q/2}L_j^{3q/2-2p}
 &\le C_{p,q,A}n^{q/4}
             \sum_{j\ge1}2^{j(3q/2-2p)}
 \le C_q n^{q/4}.
 \label{eq:dyadic-clock-summation}
\end{align}
Here $b_0\asymp\sqrt n$ and $3q/2-2p<0$. For the remaining shells,
\eqref{eq:dyadic-shell-holder} is bounded by
$C_{p,q}L_j^{q/2}e^{-c_qL_j}$, which is summable and contributes no
more than $C_qn^{q/4}$. Taking the $q$th root proves
\eqref{eq:quarter-stopping}. There is no logarithmic count of shells:
the shell contributions form a convergent geometric series. For $q=2$
one can take $p=2$, in which case the moderate-shell sum is
$C\sum n/L_j\le C\sqrt n$.

Finally, recenter at the actual mark and use
$|e^{ix}-e^{iy}|\le|x-y|$. Since $a$ vanishes outside the section,
$|a|\nu\le c_*M_a\nu_R^*$ there. Equation~\eqref{eq:quarter-stopping}
with $q=1$ proves \eqref{eq:quarter-characteristic-stopping}; its right
side is exactly zero at $v=0$. No martingale representation, conditional
independence, or mixing of the induced map is used.
\end{proof}

\subsection{One anchored factor in an arbitrary fixed moment}
For an integer $d\ge1$ and a symmetric positive semidefinite matrix
$D$, write
\[
 \mathcal G_d(D)=\mathbb E\,\mathcal Z_D^{\otimes d},\qquad
 \mathcal Z_D\sim N(0,D),\qquad
 \rho_d=\begin{cases}1/2,&d\text{ odd},\\1,&d\text{ even}.\end{cases}
\]
The tensor is zero for odd $d$. Its norm below may be any fixed
Euclidean tensor norm.

\begin{lemma}[A fixed marked collision moment]
\label{lem:all-anchored-collision-moments}
Let $a$ be complex and bounded-BV on the collision cylinder, with
$\alpha=\int a\dd\nu$. For a deterministic two-sided interval of
length $m=r+s\ge1$,
\begin{align}
 \left\|\int_M a\,\mathsf H_{r,s,R}^{\otimes d}\dd\nu
                  -\alpha m^{d/2}\mathcal G_d(\Gamma_R)\right\|
 \le C_d(\|a\|_\infty+\|a\|_{\mathrm{BV}})m^{d/2-\rho_d}.
 \label{eq:all-anchored-collision-moments}
\end{align}
The constant is independent of the position of zero in the interval.
\end{lemma}
\begin{proof}
First take a real unit direction $\xi$ and put $u=\xi\cdot h_R$.
Consider a joint cumulant containing $a-\alpha$ at time zero and
$j\ge1$ copies of $u$ at integer times $i_1,\ldots,i_j$.
The independent-block comparison in
Proposition~\ref{prop:connected-collision-summability} applies also
to these different observables: the finite-order decoupling estimate
is multilinear, and the partition polynomial is a finite sum of
products of subset moments. After normalization of $a-\alpha$, its
absolute value is bounded by
\[
 C_j(\|a\|_\infty+\|a\|_{\mathrm{BV}})
    e^{-c_j\operatorname{diam}\{0,i_1,\ldots,i_j\}}.
\]
It is summable over all $\Z^j$. Thus the cumulant containing the
single centered mark and $j$ copies of the sum of $u$ on any interval
is bounded by $C_j(\|a\|_\infty+\|a\|_{\mathrm{BV}})$, independently
of its length or position. Complex marks are handled by real and
imaginary parts.

Expand the moment with one centered mark by the finite
moment--cumulant partition identity. The block containing that mark
must contain at least one copy of the sum. Every other nonzero block
contains at least two copies, because $u$ is centered. At most
$\lfloor(d-1)/2\rfloor$ such blocks remain. Each pure sum cumulant
is $O_d(m)$ by the inherited connected-correlation theorem. Hence
\[
 \left|\int(a-\alpha)(\xi\cdot\mathsf H_{r,s,R})^d\dd\nu\right|
 \le C_d(\|a\|_\infty+\|a\|_{\mathrm{BV}})
                         m^{\lfloor(d-1)/2\rfloor}.
\]
For the unweighted moment, all cumulants of order at least two are
$O_d(m)$ and the second is $m\xi^{\mathsf T}\Gamma_R\xi+O(1)$.
If $d$ is even, the pair partitions give the Gaussian term; replacing
one covariance by its bounded correction, or using a nonpair
partition, lowers the power of $m$ by at least one. If $d$ is odd,
there are at most $(d-1)/2$ nonzero blocks and the Gaussian moment
is zero. For $d=1$ the unweighted moment is exactly zero.
These observations prove \eqref{eq:all-anchored-collision-moments}
against every $\xi^{\otimes d}$. Polarization in the fixed dimension
four gives the symmetric tensor assertion. All cumulant orders used
are fixed, at most $d+1$.
\end{proof}

\begin{theorem}[All fixed Gaussian moments with an actual return mark]
\label{thm:all-marked-gaussian-moments}
For every fixed integer $d\ge1$, all $n\ge2$ and $0\le k\le n$,
\begin{align}
 &\left\|\int_{Y_R^*}a((F_R^*)^kx)
        \left(\frac{U_{n,R}(x)}{\sqrt n}\right)^{\otimes d}\dd\nu(x)
                   -\alpha_a\mathcal G_d(D_R)\right\|\notag\\
 &\hspace{18mm}\le C_d\left(M_an^{-1/4}
                           +(M_a+V_a)n^{-\rho_d}\right).
 \label{eq:all-marked-gaussian-moments}
\end{align}
The convention is the restricted collision measure, with
$\alpha_a=\int a\dd\nu$; no positivity of $a$ is required.
In particular the actual normalized covariance and its induced
Cesaro Green--Kubo expression converge to $D_R$ at rate $O(n^{-1/4})$.
\end{theorem}
\begin{proof}
Recenter at $y=(F_R^*)^kx$. The original numerator is
\[
 \int_{Y_R^*}a(y)
       \left(\frac{Z_{n,k,R}(y)}{\sqrt n}\right)^{\otimes d}\dd\nu(y).
\]
For $d\ge2$, the tensor telescoping inequality
\[
 \|x^{\otimes d}-y^{\otimes d}\|
       \le C_d|x-y|(|x|+|y|)^{d-1}
\]
and H\"older's inequality with exponents $d$ and $d/(d-1)$ show
that its change on replacing $Z$ by $W$ is at most $C_dM_an^{-1/4}$,
by \eqref{eq:all-stopped-moment-bound} and \eqref{eq:quarter-stopping}.
For $d=1$ use the first moment version directly. Restriction from
$\nu$ to the section changes only the fixed factor $c_*$.

Apply Lemma~\ref{lem:all-anchored-collision-moments} to $W$ with
$m=m_k+m_{n-k}=n/c_*+O(1)$. Division by $n^{d/2}$ gives the second
error in \eqref{eq:all-marked-gaussian-moments}. The Gaussian term
is $\alpha_a\mathcal G_d((m/n)\Gamma_R)$, which differs from the
stated term by $O_d(M_a/n)$ for even $d$ and is zero for odd $d$.
The covariance matrices are uniformly bounded. This proves the theorem.

For the unweighted assertion take $a=s_R$, so $\alpha_a=1$ and both
insertion norms are uniformly bounded. The Kac mean makes the actual
first centered moment zero. At $d=2$ the theorem is therefore a
covariance estimate. Expanding that exact covariance under stationarity
of $F_R^*$ gives the Cesaro expression in
\eqref{eq:induced-cesaro-covariance}, with the improved error.
Absolute summability of induced correlations is not asserted.
\end{proof}

\begin{corollary}[Polynomial moments under the unchanged event]
\label{cor:all-same-event-moments}
Let $E_{n,R}\subset Y_R^*$ have probability $p_{n,R}>0$ and
zero-extended indicator variation $V_{n,R}$. For
$A_{n,k,R}=\{x:(F_R^*)^kx\in E_{n,R}\}$ and every fixed $d\ge1$,
\begin{align}
 &\left\|\mathbb E_{\nu_R^*}
   \left[\left(U_{n,R}/\sqrt n\right)^{\otimes d}\mid A_{n,k,R}\right]
                       -\mathcal G_d(D_R)\right\|\notag\\
 &\hspace{18mm}\le\frac{C_d}{p_{n,R}}
         \left(n^{-1/4}+(1+V_{n,R})n^{-\rho_d}\right).
 \label{eq:all-same-event-moments}
\end{align}
For $p_{n,R}\ge cn^{-\beta}$ and $V_{n,R}\le Cn^\kappa$,
$\beta<1/4$ and $\beta+\kappa<\rho_d$ suffice. Every fixed
polynomial moment and the conditional covariance converge under
$\beta<1/4$ and $\beta+\kappa<1/2$.
\end{corollary}
\begin{proof}
Take $a=c_*^{-1}\one_{E_{n,R}}$ in
Theorem~\ref{thm:all-marked-gaussian-moments}. Its mass equals
$p_{n,R}$, also the probability of $A_{n,k,R}$ by invariance.
Divide by this unchanged denominator. The conditional first moment
tends to zero under the stronger stated conditions, so subtracting
its tensor square from the conditional second moment proves the
covariance assertion. A fixed polynomial is a finite linear
combination of the displayed tensor coordinates.
\end{proof}

\paragraph{Scope of the moment improvement.}
The fourth-root rate is obtained for the actual unbounded record at
one prescribed count and mark. It is not a maximal-in-$k$ theorem,
a statement about growing moment orders, or an event-replacement
estimate. In particular it neither assumes induced mixing nor treats
a first-variation bound as a bound for inverse-coarea second derivatives.
